% =============================================================================
% Project : THE PHILOSOPHY OF VIABILITY
% Book    : Book I — Ontology of Viability
% Book ID : VS.B01
% Part    : Part I — Identity and Manifestation
% Part ID : VS.B01.P01
% Chapter : Chapter 08 — Comparative Philosophical Examination
% Chapter ID : VS.B01.P01.C08
%
% File :
%   4_ASSETS/tables/
%   VS.B01.P01.C08.S10.SS14.P003.TL01.kant_claims_and_ontological_limits.tex
%
% Section :
%   Section 10 — Immanuel Kant
%
% Asset :
%   Table TL01 — Kant's Claims and Their Ontological Limits
%
% Structural Identifier :
%   VS.B01.P01.C08.S10.SS14.P003.TL01
%
% Asset Type : TL
%
% Purpose :
%   To distinguish the principal claims in Kant's critical
%   philosophy according to their function, explanatory
%   reach, and limits for an ontological examination of
%   identity.
%
% Intellectual Benefit :
%   The table prevents the conflation of apperceptive unity,
%   empirical self-consciousness, substantiality, permanence,
%   and personal identity.
%
% Aesthetic / Experiential Benefit :
%   It provides a compact comparative overview after the
%   extended argument and reduces the cognitive burden of
%   retaining several closely related Kantian distinctions.
%
% Representational Limit :
%   The table is an analytical summary. It does not imply
%   that Kant's critical philosophy is internally uncontested
%   or that the comparative limits listed here settle every
%   interpretive dispute.
%
% =============================================================================

\begin{center}
\small
\renewcommand{\arraystretch}{1.18}
\setlength{\tabcolsep}{4pt}

\begin{longtable}{
    >{\raggedright\arraybackslash}p{2.3cm}
    >{\raggedright\arraybackslash}p{0.27\textwidth}
    >{\raggedright\arraybackslash}p{0.27\textwidth}
    >{\raggedright\arraybackslash}p{0.27\textwidth}
}

\caption{Kant's Principal Claims and Their Ontological Limits}
\label{tab:VS.B01.P01.C08.S10.SS14.P003.TL01} \\

\toprule
\textbf{Problem}
&
\textbf{Kant's account}
&
\textbf{What it establishes}
&
\textbf{What it does not establish}
\\
\midrule
\endfirsthead

\toprule
\textbf{Problem}
&
\textbf{Kant's account}
&
\textbf{What it establishes}
&
\textbf{What it does not establish}
\\
\midrule
\endhead

\midrule
\multicolumn{4}{r}{\footnotesize\textit{Continued on next page}}\\
\endfoot

\bottomrule
\endlastfoot

Copernican turn
&
Objects, insofar as they are objects of possible
experience, conform to a priori conditions of
human cognition.
&
The transcendental conditions under which
objective experience is possible.
&
That the existence or ontological constitution
of beings is produced by the subject.
\\

\addlinespace

Categories
&
Pure concepts of the understanding provide
rules for synthesizing the manifold of intuition.
&
The conceptual conditions under which objects
can be thought in possible experience.
&
That the categories constitute an exhaustive
inventory of the determinations of Being.
\\

\addlinespace

Transcendental apperception
&
The manifold must be capable of belonging
to one self-consciousness under the unity
of the ``I think.''
&
The formal unity required for the synthesis
of representations in one experience.
&
That the subject is thereby known as a
simple or enduring substance.
\\

\addlinespace

Transcendental deduction
&
The categories are justified through their
necessary relation to the possibility of
experience.
&
The objective validity of the categories
within their legitimate domain.
&
Their unrestricted application beyond
possible experience.
\\

\addlinespace

First Analogy
&
Substance is permanent in all change
of appearances.
&
A principle for the objective temporal
determination of appearances.
&
That every particular being or person
is thereby proved numerically identical
through all transformations.
\\

\addlinespace

Paralogisms
&
The formal unity of the thinking subject
does not prove the substantiality,
simplicity, or persistence of the soul.
&
The limits of rational psychology's
inferences from the ``I think.''
&
That the subject does not exist or
that personal identity is merely
a bundle of perceptions.
\\

\addlinespace

Personal identity
&
The identity of the ``I'' in thought
cannot alone establish the identity
of a person as a persisting substance.
&
A distinction between apperceptive
identity and substantial personal
persistence.
&
A complete positive ontology of
personal identity.
\\

\addlinespace

Phenomena and noumena
&
The categories have objective use
only under the conditions of possible
experience.
&
A critical boundary on theoretical
knowledge.
&
A determinate description of things
in themselves as objects of
theoretical knowledge.
\\

\addlinespace

Identity and difference
&
Objective comparison requires
transcendental reflection on the
relation of representations to
the faculties of cognition.
&
The insufficiency of merely logical
comparison for every judgment
about objects.
&
A complete account of the ontological
ground of identity and distinction.
\\

\end{longtable}

\vspace{0.5em}

\begin{minipage}{0.95\textwidth}
\footnotesize
\textit{Note:}
The table distinguishes the scope of Kant's transcendental
arguments from the broader ontological conclusions that
might be drawn from them. Such conclusions require
additional argument and must not be treated as automatic
consequences of the critical philosophy.
\end{minipage}

\end{center}

% =============================================================================
% END OF TABLE TL01
% =============================================================================